% Dejenerasyon: tek skaler gözlem, parametre uzayında bir NOKTA değil bir
% KORİDOR seçer.
%
% Tek eksen. Gri eş-beta düzey eğrileri, gözlemin 1-sigma koridoru ve o
% koridorun yerel doğrusallaştırması olan Fisher elipsi. Elipsin uzun ekseni
% gözlemin kısıtlamadığı yön, kısa ekseni kısıtladığı yöndür -- projenin
% "n_yön" ölçütü tam olarak bunu sayar.
%
% Eğrilerin KONUMU şematiktir (ölçüm değil): elimizde yalnız sabit
% (alpha_b, f) noktasında üç bağdaşım koşusu var. Gerekçe
% scripts/rapor2204_veri.py içinde; altyazıda da yazılı.
%
% `docs/rapor2204` içinden derle:  lualatex sekil/dejenerasyon.tex
\documentclass[border=3pt,tikz]{standalone}
\input{sekil/ortak.tex}
\usepgfplotslibrary{fillbetween}

\colorlet{cizgi}{black!78}
\colorlet{duzey}{black!32}
\colorlet{gozlem}{vurgu}

\pgfplotstableread{veri/dejenerasyon_egri.dat}\egri

% Düzey eğrisinin üstünde boşluk açıp sayıyı oraya koyar: haritacılık
% deyimi. Etiketi eğrinin yanına iliştirmekten çok daha temiz okunur.
\newcommand{\duzeyetiket}[3]{%
  \node[fill=white, inner sep=1.1pt, scale=0.78, text=duzey!85!black,
        rotate=#3] at (axis cs:#1) {#2};}

\begin{document}
\begin{tikzpicture}

\begin{axis}[
  width=82mm, height=58mm, scale only axis,
  ymode=log,
  xmin=1.0, xmax=1.30, ymin=1, ymax=1e5,
  xlabel={blok gözenekliliği $\alpha_b$},
  ylabel={matris bağdaşım dayanımı $Y_0$\ [Pa]},
  xtick={1.00,1.05,1.10,1.15,1.20,1.25,1.30},
  xticklabels={$1{,}00$,,$1{,}10$,,$1{,}20$,,$1{,}30$},
  ytick={1,1e1,1e2,1e3,1e4,1e5},
  yticklabels={$1$,,$10^{2}$,,$10^{4}$,},
  minor tick num=0,
  axis line style={draw=cizgi, line width=0.4pt},
  tick style={draw=cizgi, line width=0.4pt},
  tick align=outside,
  ticklabel style={font=\footnotesize, text=cizgi},
  label style={font=\footnotesize, text=black},
  xlabel style={yshift=1pt}, ylabel style={yshift=-2pt},
  axis on top=false,
  clip=true,
]

  % --- gözlemin 1-sigma koridoru -------------------------------------
  % Gözlemle uyumlu kümenin kendisi. Bir nokta değil; eksen boyunca uzanan
  % dar bir şerit.
  \addplot[draw=none, name path=korA, forget plot]
    table[x=alphab, y=kor_ust] {\egri};
  \addplot[draw=none, name path=korB, forget plot]
    table[x=alphab, y=kor_alt] {\egri};
  \addplot[gozlem!12, forget plot] fill between[of=korA and korB];

  % --- komşu beta düzeyleri: ince gri --------------------------------
  \addplot[duzey, line width=0.4pt, forget plot]
    table[x=alphab, y=b20] {\egri};
  \addplot[duzey, line width=0.4pt, forget plot]
    table[x=alphab, y=b25] {\egri};
  \addplot[duzey, line width=0.4pt, forget plot]
    table[x=alphab, y=b30] {\egri};
  \addplot[duzey, line width=0.4pt, forget plot]
    table[x=alphab, y=b40] {\egri};
  \addplot[duzey, line width=0.4pt, forget plot]
    table[x=alphab, y=b45] {\egri};

  % --- gözlenen beta düzeyi ------------------------------------------
  \addplot[gozlem, line width=0.9pt, forget plot]
    table[x=alphab, y=bgoz] {\egri};

  % --- düzey etiketleri, eğrinin İÇİNDE ------------------------------
  % Açı sayfa ölçeğinden: 0,30 birim -> 82 mm, 5 ondalık -> 58 mm;
  % eğim -12 ondalık/birim => atan(-12*(58/5)/(82/0,30)) = -27 derece.
  %
  % Konumlar elle değil, eğrinin kendi denkleminden:
  %   log10 Y0 = L(beta) - 12 (alpha_b - 1),  L = 3,86 - 1,2 (beta - 3,5418)
  % Yatay yerler, etiketler bir sütun oluşturmasın diye kaydırıldı.
  \duzeyetiket{1.085,4.90e4}{$\beta=2{,}0$}{-27}   % L=5,710
  \duzeyetiket{1.155,1.78e3}{$2{,}5$}{-27}         % L=5,110
  \duzeyetiket{1.225,6.46e1}{$3{,}0$}{-27}         % L=4,510
  \duzeyetiket{1.065,3.39e2}{$4{,}0$}{-27}         % L=3,310
  \duzeyetiket{1.020,2.95e2}{$4{,}5$}{-27}         % L=2,710

  \coordinate (fis) at (axis cs:1.150,1.148e2);
  \coordinate (ucA) at (axis cs:1.016,4.60e3);

\end{axis}

% ===== Fisher elipsi: sayfa düzleminde, koridorla aynı açıda ==========
% Dolgu YOK: elips koridorun yerel doğrusallaştırmasıdır, onun üstüne
% konan ayrı bir nesne değil. Beyaz dolgu bu ilişkiyi gizliyordu.
% Uzun yarı eksen 9,5 mm: yerel bir yaklaşım olduğu görünsün, koridoru
% baştan sona kaplamasın.
\begin{scope}[shift={(fis)}, rotate=-27]
  \draw[cizgi, line width=0.6pt]
       (0,0) ellipse [x radius=9.5mm, y radius=2.4mm];
  \draw[cizgi, line width=0.4pt] (-9.5mm,0) -- (9.5mm,0);
  \draw[cizgi, line width=0.4pt] (0,-2.4mm) -- (0,2.4mm);
\end{scope}

% ===== açıklamalar: kısa, boş alana, ince kılavuzla ===================
% Uzun açıklama altyazıya bırakıldı; şekil üstünde yalnız ad kalır.
% Kılavuzlar elipsin uçlarından başlar: uzun eksen ucu -27 derecede
% (8,5; -4,3) mm, kısa eksen ucu buna dik (1,1; 2,1) mm.
\draw[cizgi, line width=0.3pt]
     ($(fis)+(8.5mm,-4.3mm)$) -- ($(fis)+(13.0mm,-8.8mm)$);
\node[anchor=north west, font=\footnotesize, text=black, inner sep=1.5pt,
      fill=white] at ($(fis)+(12.6mm,-8.4mm)$) {kısıtlanmayan yön};

\draw[cizgi, line width=0.3pt]
     ($(fis)+(1.1mm,2.1mm)$) -- ($(fis)+(5.6mm,7.6mm)$);
\node[anchor=south west, font=\footnotesize, text=black, inner sep=1.5pt,
      fill=white] at ($(fis)+(5.2mm,7.3mm)$) {kısıtlanan yön};

\draw[gozlem!75, line width=0.3pt]
     ($(ucA)+(1.2mm,1.4mm)$) -- ($(ucA)+(6.0mm,6.2mm)$);
\node[anchor=south west, align=left, font=\footnotesize,
      text=gozlem!85!black, inner sep=1pt]
     at ($(ucA)+(5.6mm,5.9mm)$)
     {$\beta_{\text{göz}}$ ile uyumlu küme};

\end{tikzpicture}
\end{document}
